\documentclass{article}

\usepackage[T1]{fontenc}

\usepackage{booktabs,graphicx,siunitx}

\begin{document}

\section{One-sample t-test}

\begin{tabular}{lllll}
\toprule
Variable & N & M & SD & SE \\
\midrule
pretest\_score & 240 & 49.25 & 8.46 & 0.55 \\
\bottomrule
\end{tabular}

\begin{tabular}{lllllll}
\toprule
Test value & t & df & p & Mdiff & 95\% CI & Cohen's d [95\% CI] \\
\midrule
50 & $-$1.38 & 239 & .170 & $-$0.75 & [$-$1.83, 0.32] & $-$0.09 [$-$0.22, 0.04] \\
\bottomrule
\end{tabular}

\includegraphics{figures/01_one-sample-t-test/figure_distribution.png}

Compares your sample mean to a fixed value. A p below alpha means the mean differs from that value; the mean difference and 95\% CI show by how much (its sign tells you whether the mean is above or below the test value), and Cohen's d gives the effect size. A p at or above alpha does not prove the mean equals the test value --- only that you lack evidence it differs. Your significance threshold ($\alpha$) is 0.05.

A one-sample t-test gave M=49.2 vs. 50, t(239)=$-$1.38, p = .170, d=$-$0.09 [$-$0.22, 0.04].

Statistical basis: One-sample t-test (Student (1908). "The probable error of a mean." Biometrika, 6(1), 1-25.); Cohen's d effect-size benchmarks (.2/.5/.8) (Cohen, J. (1988). Statistical Power Analysis for the Behavioral Sciences (2nd ed.). Routledge.).

\end{document}